\chapter*{Abkürzungsverzeichnis}
\thispagestyle{plain}

% Alphabetisch sortiert (Groß-/Kleinschreibung ignoriert).
% Aufgenommen sind die fachspezifischen Abkürzungen, die in der Arbeit
% mehrfach verwendet werden. Bewusst nicht aufgenommen:
% - Eigennamen von Datensätzen, Konferenzen, Software und Institutionen
%   (MultiNERD, CrossNER, OntoNotes, SemEval, CUDA, SLURM, DFKI, ...);
%   Modellnamen nur, wenn sie selbst Akronyme sind (BERT, GPT, ...),
% - allgemein bekannte IT-Abkürzungen (API, JSON, XML, YAML, CPU, ...),
% - Abkürzungen, die nur an einer Stelle vorkommen und dort im Text
%   erklärt werden (NSP, DoRA, rsLoRA, KV, ...),
% - seltene Tagging-Schema-Varianten (IOBES, BILOU, IOE), die in
%   Kapitel 2 erläutert werden; das Tagging-Schema IOB2 (auch BIO) steht
%   unter IOB, da BIO in MultiNERD zugleich ein Entitätstyp ist,
% - Maßeinheiten (GB, MiB, ms) und Entitätstypen-Labels (PER, LOC, ...).

\begin{singlespace}
\begin{acronym}[RoBERTa]
    \setlength{\parskip}{3pt}  % Zeilenabstand 1,5 plus etwas Luft zwischen den Einträgen
    \acro{BERT}{Bidirectional Encoder Representations from Transformers}
    \acro{BF16}{BFloat16 (Brain Floating Point, 16 Bit)}
    \acro{BiLSTM}{Bidirectional Long Short-Term Memory}
    \acro{CoNLL}{Conference on Computational Natural Language Learning}
    \acro{CRF}{Conditional Random Field}
    \acro{DeBERTa}{Decoding-enhanced BERT with Disentangled Attention}
    \acro{ELECTRA}{Efficiently Learning an Encoder that Classifies Token Replacements Accurately}
    \acro{F1}[F$_1$]{F$_1$-Score, harmonisches Mittel aus Precision und Recall}
    \acro{FF}{Forschungsfrage}
    \acro{FN}{False Negatives}
    \acro{FP}{False Positives}
    \acro{FP32}{32-Bit-Gleitkommaformat (Single Precision)}
    \acro{GPT}{Generative Pre-trained Transformer}
    \acro{GPU}{Graphics Processing Unit}
    \acro{HBM3}{High Bandwidth Memory 3 (GPU-Speichertechnologie)}
    \acro{IE}{Information Extraction}
    \acro{IOB}{Inside, Outside, Beginning (Tagging-Schema; Varianten IOB1 und IOB2, auch BIO genannt)}
    \acro{LLaMA}{Large Language Model Meta AI}
    \acro{LLM}{Large Language Model}
    \acro{LoRA}{Low-Rank Adaptation}
    \acro{MLM}{Masked Language Modeling}
    \acro{MRC}{Machine Reading Comprehension}
    \acro{MTP}{Multi-Token-Prediction}
    \acro{NER}{Named Entity Recognition}
    \acro{NF4}{4-Bit-NormalFloat}
    \acro{NLP}{Natural Language Processing}
    \acro{PEFT}{Parameter-Efficient Fine-Tuning}
    \acro{QLoRA}{Quantized Low-Rank Adaptation}
    \acro{RE}{Relation Extraction}
    \acro{RoBERTa}{Robustly Optimized BERT Pre-training Approach}
    \acro{SFT}{Supervised Fine-Tuning}
    \acro{TP}{True Positives}
    \acro{UIE}{Universal Information Extraction}
\end{acronym}
\end{singlespace}
